\answerpaper{非数学类}

\answerproblem{一、}{30}填空题（每小题 6 分）：
\begin{enumerate}[label=(\arabic*)]
  \item $\displaystyle\lim_{x\to+\infty}\sqrt{x^{2}+x+1}\cdot\frac{x-\ln(\ee^{x}+x)}{x}$；
  \item 设 $z=z(x,y)$ 是由方程 $2\sin(x+2y-3z)=x+2y-3z$ 所确定的二元隐函数，求 $\dfrac{\partial z}{\partial x}+\dfrac{\partial z}{\partial y}$；
  \item 设函数 $f(x)$ 连续，且 $f(0)\neq0$，求 $\displaystyle\lim_{x\to0}\frac{2\int_{0}^{x}(x-t)f(t)\dd t}{x\int_{0}^{x}f(x-t)\dd t}$；
  \item 求过三条直线
  \[
  L_{1}:\begin{cases}x=0,\\ y-z=2,\end{cases}\qquad
  L_{2}:\begin{cases}x=0,\\ x+y-z+2=0,\end{cases}\qquad\text{与}\qquad
  L_{3}:\begin{cases}x=\sqrt{2},\\ y-z=0\end{cases}
  \]
  的圆柱面方程；
  \item 记 $D=\left\{(x,y)\mid x^{2}+y^{2}\le\pi\right\}$，求 $\displaystyle\iint_{D}\left(\sin x^{2}\cos y^{2}+x\sqrt{x^{2}+y^{2}}\right)\dd x\dd y$．
\end{enumerate}
\begin{solution}
\textbf{(1)} 原式
\[
=\lim_{x\to+\infty}\sqrt{x^{2}+x+1}\cdot\frac{x-\ln(\ee^{x}+x)}{x}
=-\lim_{x\to+\infty}\sqrt{1+\frac{1}{x}+\frac{1}{x^{2}}}\ln\left(1+\frac{x}{\ee^{x}}\right)=0.
\]

\textbf{(2)} 将方程两边分别关于 $x$ 和 $y$ 求偏导，得
\[
\begin{cases}
2\cos(x+2y-3z)\left(1-3\dfrac{\partial z}{\partial x}\right)=1-3\dfrac{\partial z}{\partial x},\\[2ex]
2\cos(x+2y-3z)\left(2-3\dfrac{\partial z}{\partial y}\right)=2-3\dfrac{\partial z}{\partial y}.
\end{cases}
\]
按 $\cos(x+2y-3z)=\dfrac{1}{2}$ 和 $\neq\dfrac{1}{2}$ 两种情形，都可解得
\[
\begin{cases}
\dfrac{\partial z}{\partial x}=\dfrac{1}{3},\\[2ex]
\dfrac{\partial z}{\partial y}=\dfrac{2}{3},
\end{cases}
\]
因此 $\dfrac{\partial z}{\partial x}+\dfrac{\partial z}{\partial y}=1$．

\textbf{(3)} 原式
\begin{align*}
&=\lim_{x\to0}\frac{2x\int_{0}^{x}f(t)\dd t-2\int_{0}^{x}tf(t)\dd t}{x\int_{0}^{x}f(u)\dd u}
=\lim_{x\to0}\frac{2\int_{0}^{x}f(t)\dd t+2xf(x)-2xf(x)}{\int_{0}^{x}f(u)\dd u+xf(x)}\\
&=\lim_{x\to0}\frac{2\int_{0}^{x}f(t)\dd t}{\int_{0}^{x}f(u)\dd u+xf(x)}
=\lim_{x\to0}\frac{2xf(\xi)}{xf(\xi)+xf(x)}=1,
\end{align*}
其中 $\xi$ 介于 $0$，$x$ 之间．

\textbf{(4)} 三条直线的对称式方程分别为
\[
L_{1}:\frac{x}{0}=\frac{y-1}{1}=\frac{z+1}{1},\qquad
L_{2}:\frac{x}{0}=\frac{y-0}{1}=\frac{z-2}{1},\qquad
L_{3}:\frac{x-\sqrt{2}}{0}=\frac{y-1}{1}=\frac{z-1}{1},
\]
所以三条直线平行．

在 $L_{1}$ 上取点 $P_{1}(0,1,-1)$，过该点作与三直线都垂直的平面 $y+z=0$，分别交 $L_{2}$，$L_{3}$ 于点 $P_{2}(0,-1,1)$，$P_{3}(\sqrt{2},0,0)$．易知经过这三点的圆的圆心为 $O(0,0,0)$．这样，所求圆柱面的中心轴线方程为 $\dfrac{x}{0}=\dfrac{y}{1}=\dfrac{z}{1}$．

设圆柱面上任意点的坐标为 $Q(x,y,z)$，因为点 $Q$ 到轴线的距离均为 $\sqrt{2}$，所以有
\[
\frac{\abs{(x,y,z)\times(0,1,1)}}{\sqrt{0^{2}+1^{2}+1^{2}}}=\sqrt{2},
\]
化简即得所求圆柱面的方程为 $2x^{2}+y^{2}+z^{2}-2yz=4$．

\textbf{(5)} 根据重积分的对称性，得
\begin{align*}
\text{原式}&=\iint_{D}\sin x^{2}\cos y^{2}\dd x\dd y=\iint_{D}\sin y^{2}\cos x^{2}\dd x\dd y\\
&=\frac{1}{2}\iint_{D}\left(\sin x^{2}\cos y^{2}+\sin y^{2}\cos x^{2}\right)\dd x\dd y\\
&=\frac{1}{2}\iint_{D}\sin(x^{2}+y^{2})\dd x\dd y=\frac{1}{2}\int_{0}^{2\pi}\dd\theta\int_{0}^{\sqrt{\pi}}r\sin r^{2}\dd r\\
&=\frac{\pi}{2}\left(-\cos r^{2}\right)\Big|_{0}^{\sqrt{\pi}}=\pi.
\end{align*}
\end{solution}

\answerproblem{二、}{14}设 $x_{1}=2021$，$x_{n}^{2}-2(x_{n}+1)x_{n+1}+2021=0$（$n\ge1$），证明数列 $\{x_{n}\}$ 收敛，并求极限 $\lim\limits_{n\to\infty}x_{n}$．
\begin{solution}
记 $a=1011$，$y_{n}=1+x_{n}$，函数 $f(x)=\dfrac{x}{2}+\dfrac{a}{x}$（$x>0$），则 $y_{1}=2a$，且
\[
y_{n+1}=f(y_{n})\quad(n\ge1).
\]
易知，当 $x>\sqrt{2a}$ 时，$x>f(x)>\sqrt{2a}$，所以 $\{y_{n}\}$ 是单调减少且有下界的数列，因而收敛．由此可知 $\{x_{n}\}$ 收敛．

令 $\lim\limits_{n\to\infty}y_{n}=A$，则 $A>0$ 且 $A=f(A)$，解得 $A=\sqrt{2a}$．因此 $\lim\limits_{n\to\infty}x_{n}=\sqrt{2022}-1$．
\end{solution}

\answerproblem{三、}{14}设 $f(x)$ 在 $[0,+\infty)$ 上是有界连续函数，证明：方程 $y''+14y'+13y=f(x)$ 的每一个解在 $[0,+\infty)$ 上都是有界函数．
\begin{solution}
易得对应的齐次方程 $y''+14y'+13y=0$ 的通解为
\[
y=C_{1}\ee^{-x}+C_{2}\ee^{-13x}.
\]
又由 $y''+14y'+13y=f(x)$，得 $(y''+y')+13(y'+y)=f(x)$．令 $y_{1}=y'+y$，则 $y_{1}'+13y_{1}=f(x)$，解得
\[
y_{1}=\ee^{-13x}\left(\int_{0}^{x}f(t)\ee^{13t}\dd t+C_{3}\right).
\]
同理，由 $y''+14y'+13y=f(x)$，得 $(y''+13y')+(y'+13y)=f(x)$．令 $y_{2}=y'+13y$，则 $y_{2}'+y_{2}=f(x)$，解得
\[
y_{2}=\ee^{-x}\left(\int_{0}^{x}f(t)\ee^{t}\dd t+C_{4}\right).
\]
取 $C_{3}=C_{4}=0$，得
\[
\begin{cases}
y'+y=\ee^{-13x}\int_{0}^{x}f(t)\ee^{13t}\dd t,\\[2ex]
y'+13y=\ee^{-x}\int_{0}^{x}f(t)\ee^{t}\dd t,
\end{cases}
\]
由此解得原方程的一个特解为
\[
y^{*}=\frac{1}{12}\ee^{-x}\int_{0}^{x}f(t)\ee^{t}\dd t-\frac{1}{12}\ee^{-13x}\int_{0}^{x}f(t)\ee^{13t}\dd t.
\]
因此，原方程的通解为
\[
y=C_{1}\ee^{-x}+C_{2}\ee^{-13x}+\frac{1}{12}\ee^{-x}\int_{0}^{x}f(t)\ee^{t}\dd t-\frac{1}{12}\ee^{-13x}\int_{0}^{x}f(t)\ee^{13t}\dd t.
\]
因为 $f(x)$ 在 $[0,+\infty)$ 上有界，所以存在 $M>0$，使得 $\abs{f(x)}\le M$，$0\le x<+\infty$，注意到当 $x\in[0,+\infty)$ 时，$0<\ee^{-x}\le1$，$0<\ee^{-13x}\le1$，所以
\begin{align*}
\abs{y}&\le\abs{C_{1}\ee^{-x}}+\abs{C_{2}\ee^{-13x}}+\frac{1}{12}\ee^{-x}\abs{\int_{0}^{x}f(t)\ee^{t}\dd t}+\frac{1}{12}\ee^{-13x}\abs{\int_{0}^{x}f(t)\ee^{13t}\dd t}\\
&\le\abs{C_{1}}+\abs{C_{2}}+\frac{M}{12}\ee^{-x}\int_{0}^{x}\ee^{t}\dd t+\frac{M}{12}\ee^{-13x}\int_{0}^{x}\ee^{13t}\dd t\\
&\le\abs{C_{1}}+\abs{C_{2}}+\frac{M}{12}\left(1-\ee^{-x}\right)+\frac{M}{12\times13}\left(1-\ee^{-13x}\right)\\
&\le\abs{C_{1}}+\abs{C_{2}}+\frac{M}{12}+\frac{M}{12\times13}=\abs{C_{1}}+\abs{C_{2}}+\frac{7M}{78}.
\end{align*}
对于方程的每一个确定的解，常数 $C_{1}$、$C_{2}$ 是固定的，所以，原方程的每一个解都是有界的．
\end{solution}

\answerproblem{四、}{14}对于 4 次齐次函数
\[
f(x,y,z)=a_{1}x^{4}+a_{2}y^{4}+a_{3}z^{4}+3a_{4}x^{2}y^{2}+3a_{5}y^{2}z^{2}+3a_{6}x^{2}z^{2},
\]
计算曲面积分 $\displaystyle\iint_{\Sigma}f(x,y,z)\dd S$，其中 $\Sigma$：$x^{2}+y^{2}+z^{2}=1$．
\begin{solution}
因为 $f(x,y,z)$ 为 4 次齐次函数，所以对 $\forall t\in\R$，恒有
\[
f(tx,ty,tz)=t^{4}f(x,y,z).
\]
对上式两边关于 $t$ 求导，得
\[
xf'_{x}(tx,ty,tz)+yf'_{y}(tx,ty,tz)+zf'_{z}(tx,ty,tz)=4t^{3}f(x,y,z).
\]
取 $t=1$，得
\[
xf'_{x}(x,y,z)+yf'_{y}(x,y,z)+zf'_{z}(x,y,z)=4f(x,y,z).
\]
设曲面 $\Sigma$ 上点 $(x,y,z)$ 处的外法线方向的方向余弦为 $(\cos\alpha,\cos\beta,\cos\gamma)$，则 $\cos\alpha=x$，$\cos\beta=y$，$\cos\gamma=z$，因此
\begin{align*}
\iint_{\Sigma}f(x,y,z)\dd S
&=\frac{1}{4}\iint_{\Sigma}\Big[xf'_{x}(x,y,z)+yf'_{y}(x,y,z)\\
&\qquad\qquad\qquad\qquad+zf'_{z}(x,y,z)\Big]\dd S\\
&=\frac{1}{4}\iint_{\Sigma}\left[\cos\alpha f'_{x}(x,y,z)+\cos\beta f'_{y}(x,y,z)+\cos\gamma f'_{z}(x,y,z)\right]\dd S\\
&=\frac{1}{4}\iint_{\Sigma}\left[f'_{x}(x,y,z)\dd y\dd z+f'_{y}(x,y,z)\dd z\dd x+f'_{z}(x,y,z)\dd x\dd y\right]\\
&=\frac{1}{4}\iiint_{x^{2}+y^{2}+z^{2}\le1}\left[f''_{xx}+f''_{yy}+f''_{zz}\right]\dd x\dd y\dd z\qquad(\text{利用高斯公式})\\
&=\frac{3}{2}\iiint_{x^{2}+y^{2}+z^{2}\le1}\Big[x^{2}(2a_{1}+a_{4}+a_{6})+y^{2}(2a_{2}+a_{4}+a_{5})\\
&\qquad\qquad\qquad\qquad+z^{2}(2a_{3}+a_{5}+a_{6})\Big]\dd x\dd y\dd z\quad(\text{利用轮换对称性})\\
&=\sum_{i=1}^{6}a_{i}\iiint_{x^{2}+y^{2}+z^{2}\le1}(x^{2}+y^{2}+z^{2})\dd x\dd y\dd z\\
&=\sum_{i=1}^{6}a_{i}\int_{0}^{2\pi}\dd\theta\int_{0}^{\pi}\dd\varphi\int_{0}^{1}\rho^{2}\cdot\rho^{2}\sin\varphi\dd\rho
=\frac{4\pi}{5}\sum_{i=1}^{6}a_{i}.
\end{align*}
\end{solution}

\answerproblem{五、}{14}设函数 $f(x)$ 在闭区间 $[a,b]$ 上有连续的二阶导数，证明：
\begin{align*}
\lim_{n\to\infty}n^{2}\Bigg[\int_{a}^{b}f(x)\dd x&-\frac{b-a}{n}\sum_{k=1}^{n}f\left(a+\frac{2k-1}{2n}(b-a)\right)\Bigg]\\
&=\frac{(b-a)^{2}}{24}\left[f'(b)-f'(a)\right].
\end{align*}
\begin{solution}
记 $x_{k}=a+\dfrac{k(b-a)}{n}$，$\xi_{k}=a+\dfrac{(2k-1)(b-a)}{2n}$，$k=1,2,\cdots,n$．将 $f(x)$ 在 $[x_{k-1},x_{k}]$ 上展开成泰勒公式，得
\[
f(x)=f(\xi_{k})+f'(\xi_{k})(x-\xi_{k})+\frac{f''(\eta_{k})}{2}(x-\xi_{k})^{2},
\]
其中 $x\in[x_{k-1},x_{k}]$，$\eta_{k}$ 介于 $0$ 和 $x$ 之间．于是
\begin{align*}
B_{n}&=\int_{a}^{b}f(x)\dd x-\frac{b-a}{n}\sum_{k=1}^{n}f\left(a+\frac{2k-1}{2n}(b-a)\right)
=\sum_{k=1}^{n}\int_{x_{k-1}}^{x_{k}}\left(f(x)-f(\xi_{k})\right)\dd x\\
&=\sum_{k=1}^{n}\int_{x_{k-1}}^{x_{k}}\left[f'(\xi_{k})(x-\xi_{k})+\frac{f''(\eta_{k})}{2}(x-\xi_{k})^{2}\right]\dd x
=\frac{1}{2}\sum_{k=1}^{n}\int_{x_{k-1}}^{x_{k}}f''(\eta_{k})(x-\xi_{k})^{2}\dd x.
\end{align*}
设 $f''(x)$ 在 $[x_{k-1},x_{k}]$ 上的最大值和最小值分别为 $M_{k}$，$m_{k}$，因为
\[
\int_{x_{k-1}}^{x_{k}}(x-\xi_{k})^{2}\dd x=\frac{(b-a)^{3}}{12n^{3}},
\]
所以
\[
\frac{(b-a)^{2}}{24}\sum_{k=1}^{n}m_{k}\frac{b-a}{n}\le n^{2}B_{n}\le\frac{(b-a)^{2}}{24}\sum_{k=1}^{n}M_{k}\frac{b-a}{n}.
\]
因为 $f''(x)$ 在 $[a,b]$ 上连续，所以 $f''(x)$ 在 $[a,b]$ 上可积．根据定积分 $\displaystyle\int_{a}^{b}f''(x)\dd x$ 的定义，及牛顿-莱布尼兹公式，得
\[
\lim_{n\to\infty}\sum_{k=1}^{n}m_{k}\frac{b-a}{n}=\lim_{n\to\infty}\sum_{k=1}^{n}M_{k}\frac{b-a}{n}=\int_{a}^{b}f''(x)\dd x=f'(b)-f'(a).
\]
再根据夹逼准则，得
\[
\lim_{n\to\infty}n^{2}B_{n}=\frac{(b-a)^{2}}{24}\left[f'(b)-f'(a)\right].
\]
\end{solution}

\answerproblem{六、}{14}设 $\{a_{n}\}$ 与 $\{b_{n}\}$ 均为正实数列，满足：$a_{1}=b_{1}=1$，且 $b_{n}=a_{n}b_{n-1}-2$，$n=2,3,\cdots$．又设 $\{b_{n}\}$ 为有界数列，证明级数 $\displaystyle\sum_{n=1}^{\infty}\frac{1}{a_{1}a_{2}\cdots a_{n}}$ 收敛，并求该级数的和．
\begin{solution}
首先，注意到 $a_{1}=b_{1}=1$，且 $a_{n}=\left(1+\dfrac{2}{b_{n}}\right)\dfrac{b_{n}}{b_{n-1}}$，所以当 $n\ge2$ 时，有
\[
a_{1}a_{2}\cdots a_{n}=\left(1+\frac{2}{b_{2}}\right)\left(1+\frac{2}{b_{3}}\right)\cdots\left(1+\frac{2}{b_{n}}\right)b_{n}.
\]
由于 $\{b_{n}\}$ 有界，故存在 $M>0$，使得当 $n\ge1$ 时，恒有 $0<b_{n}\le M$．因此
\[
0<\frac{b_{n}}{a_{1}a_{2}\cdots a_{n}}=\left(1+\frac{2}{b_{2}}\right)^{-1}\left(1+\frac{2}{b_{3}}\right)^{-1}\cdots\left(1+\frac{2}{b_{n}}\right)^{-1}
\le\left(1+\frac{2}{M}\right)^{-n+1}\to0,\ n\to\infty.
\]
根据夹逼准则，
\[
\lim_{n\to\infty}\frac{b_{n}}{a_{1}a_{2}\cdots a_{n}}=0.
\]
考虑级数 $\displaystyle\sum_{n=1}^{\infty}\frac{1}{a_{1}a_{2}\cdots a_{n}}$ 的部分和 $S_{n}$，当 $n\ge2$ 时，有
\begin{align*}
S_{n}&=\sum_{k=1}^{n}\frac{1}{a_{1}a_{2}\cdots a_{k}}
=\frac{1}{a_{1}}+\sum_{k=2}^{n}\frac{1}{a_{1}a_{2}\cdots a_{k}}\cdot\frac{a_{k}b_{k-1}-b_{k}}{2}\\
&=1+\frac{1}{2}\sum_{k=2}^{n}\left(\frac{b_{k-1}}{a_{1}a_{2}\cdots a_{k-1}}-\frac{b_{k}}{a_{1}a_{2}\cdots a_{k}}\right)\\
&=\frac{3}{2}-\frac{b_{n}}{2a_{1}a_{2}\cdots a_{n}},
\end{align*}
所以 $\lim\limits_{n\to\infty}S_{n}=\dfrac{3}{2}$，这就证明了级数 $\displaystyle\sum_{n=1}^{\infty}\frac{1}{a_{1}a_{2}\cdots a_{n}}$ 收敛，且其和为 $\dfrac{3}{2}$．
\end{solution}
